Over the past few years, there has been a rapid increase in the number of technical roles globally, leading to intensified competition
among candidates and increasing the volume of applicants that recruiters must screen. The resulting pressure on recruiters from this
candidate pool has placed strain on Curriculum Vitae (CV) reviews. This has led to a shift towards automated screening tools to
process large batches of candidates. However, the current state-of-the-art automated screening systems remain rigid, opaque, and dependent
on the CV as a sole source of information.

This report presents \textbf{M.E.R.I.T} (Multi-source Employment Ranking \& Insight Tool), a customisable prototype multi-source
screening system designed to modernise technical recruitment by moving beyond this single-source dependency. MERIT supplements a
candidate's CV with evidence from alternative sources such as GitHub and LinkedIn, with GitHub forming the primary evaluated source in
this work. The system computes flexible metrics that are derived directly from the Job Description (JD), and provides glass-box
explanations for its candidate rankings to allow recruiters to challenge the rankings, rather than succumb to automation bias. To the
best of our knowledge, no prior published system targets technical recruitment by combining extensible JD-driven weighting, multi-source
fusion, and explainable scoring in one auditable architecture.

The system's architecture is built around three key concepts, each of which ties into the stages of a screening pipeline. The first stage
is the ingestion stage, where MERIT introduces \emph{extensible} metric extraction and weighting, providing recruiters with the tooling
to define their own priorities for a particular job. The next stage is the scoring process, where MERIT leverages \emph{multi-source}
fusion to build an in-depth profile of the candidate, highlighting any alignments or contradictions between the candidate's self-reported
information and the evidence from the alternative sources. In the last stage, recruiters would often be met with raw numbers, with no
explanation or justification for the score. Instead, MERIT explores \emph{explainable} scoring, offering transparent audit trails and
source-level attribution to allow recruiters to understand the score, as well as the reasoning behind it.

Our evaluation highlights MERIT's capacity to recover qualified candidates that are often overlooked by CV-only screening systems,
especially where self-reported CV text understates demonstrated public work, as in hidden-gem profiles where rank shifts of up to eight
places were observed once GitHub evidence was fused. MERIT's ability to separate candidates on the basis of demonstrated
evidence is also presented through a Spearman correlation of $\rho = 0.733$ between MERIT (Full) and recruiter consensus in one
high-discrimination cohort ($n = 10$, exploratory). In addition, a small-scale pilot study supports the feasibility of the system, with
evidence of recruiters reversing their decisions when presented with fraud alerts, keyword stuffing, or stronger demonstrated evidence.

MERIT is presented in this report as an \textbf{auditable screening assistant} designed with the intent of supporting recruiters in their
decision-making process, rather than a fairness-certified or fully autonomous hiring authority.
